\section{Conditional projective log abundance}\label{proj:sec:introduction}

Appendices~\ref{proj:sec:introduction}--\ref{proj:sec:completion}
establish the conditional projective abundance theorem used in the main
proof. Its logarithmic-Iitaka premise,
Assumption~\ref{proj:asm:iitaka}, follows from the reduction in Lemma~\ref{red:log-iitaka}. We give the complete argument, including
its signed-boundary, current-theoretic, and Frobenius constructions.



The log abundance conjecture asserts that a nef log canonical divisor
on a projective log canonical pair is semiample. Its conclusion turns a
numerical positivity condition into a morphism defined by pluricanonical
sections. In the minimal model program, abundance complements the
existence of minimal models: a nef adjoint should determine the
appropriate canonical fibration. The existence and finite-generation
theorems for big klt adjoints are central foundations of this program
\cite{BCHM}. The lower-dimensional reductions relevant here are
developed in \cite{Proj-Hashizume2018,Proj-FujinoGongyoRings2014}.

These appendices prove the full rational-boundary assertion under one
explicit subadditivity assumption. The result is a positive
\emph{conditional} resolution of the log abundance conjecture.
It is not an unconditional abundance theorem.

\subsection{The assumption and the theorem}

\begin{assumption}[Logarithmic Iitaka subadditivity]\label{proj:asm:iitaka}
Let \(f:X\to Y\) be a surjective morphism with connected fibers between
smooth connected projective complex varieties. Let \(D_X,D_Y\) be
reduced effective simple normal crossing divisors, allowing zero
boundaries, such that
\[
  \Supp(f^*D_Y)\subseteq\Supp(D_X).
\]
For a very general smooth fiber \(F\), put \(D_F=D_X|_F\). Then
\[
 \kappa(X,K_X+D_X)\ge
 \kappa(F,K_F+D_F)+\kappa(Y,K_Y+D_Y).
\]
Here \(D_F\) is reduced with simple normal crossings and
\(K_F+D_F\sim(K_X+D_X)|_F\). The convention
\((-\infty)+b=-\infty\) applies also when \(b=-\infty\), and the zero
divisor on a point has Iitaka dimension zero.
\end{assumption}

Assumption~\ref{proj:asm:iitaka} requires no nefness, abundance, bigness,
or good-model hypothesis, and does not require \(f\) to be smooth
away from the boundary. The divisor \(D_X\) may contain additional
components. No fractional-boundary,
orbifold, nonprojective, Hodge-conjectural, or arithmetic extension is
assumed. In fact, the proof uses only \(D_X=D_Y=0\), in the Albanese
reduction in Appendix~\ref{proj:sec:exclusions}.

\begin{theorem}[Conditional log abundance]\label{proj:thm:main}\label{proj:main}\label{proj:conditional-log-abundance}
Assume Assumption~\ref{proj:asm:iitaka}. Let \(k\) be an algebraically closed
field of characteristic zero, and let \((X,B)\) be a normal projective log
canonical pair over \(k\), where \(B\) is an effective \(\Q\)-divisor
and \(K_X+B\) is \(\Q\)-Cartier. If \(K_X+B\) is nef, then it is
semiample: for some integer \(m>0\), the divisor \(m(K_X+B)\) is
Cartier and \(\OO_X(m(K_X+B))\) is generated by its global sections.
\end{theorem}

\subsection{Main constructions}

The proof proceeds by simultaneous induction on smooth nonvanishing
and existence of good log minimal models, as set out in
Appendix~\ref{proj:sec:induction}. Its principal constructions are as follows.

First, Appendices~\ref{proj:sg:section}--\ref{proj:sg:descent} prove an abundance
criterion for a nef reduced dlt adjoint
that has a \emph{signed} rational representative supported on its
boundary. Semiampleness on that boundary is supplied inductively.
A normalized root construction separates its positive and negative
parts. An infinitesimal obstruction is placed in the lowest piece of
a projective Hodge-module direct image. Negative-ample vanishing on
a finite cover of a root gerbe kills the obstruction. Algebraization
then produces compact numerically trivial subvarieties, and nef
reduction yields rational-linear descent to a big divisor on a base.

Second, Appendix~\ref{proj:sec:exclusions} subjects a hypothetical smooth
nonvanishing counterexample to geometric exclusions. Algebraic webs and positive currents force
very-general proper subvarieties to be of general type. An ascending
chain condition for the Lelong numbers of a fixed current at valuations
of bounded discrepancy converts asymptotically small intersection gaps
into exact zero gaps. This rules out canonical-class positive currents
that are pulled back from a lower-dimensional base on a dense open,
as well as moving curves of bounded normalized degree and genus.

Third, Appendices~\ref{proj:sec:jets} and~\ref{proj:fr:section} compare two jet
estimates. Moving base-locus estimates on a projective bundle over a
product create a large Seshadri constant while the corresponding
one-factor section orders remain small. After all geometric choices
have been fixed, reduction modulo large primes converts these two
estimates into incompatible determinant multiplicities. The Frobenius
diagonal filtration and a uniform curve-slope bound are the numerical
ingredients of this final comparison.

These constructions are proved below; they are not additional parts
of Assumption~\ref{proj:asm:iitaka}. Established results from the minimal
model program, Hodge modules, and positivity theory are invoked with
their hypotheses at the point of use. The order of the induction is
important: the proof of smooth nonvanishing in dimension \(n\) uses
good models only in dimensions below \(n\).
Appendix~\ref{proj:sec:completion} closes the induction and transfers the
complex conclusion to arbitrary algebraically closed fields of
characteristic zero.

\subsection{Conventions}

A variety is integral. Canonical divisors are chosen compatibly on
birational models. For a divisor \(E\) on a smooth model
\(p:W\to X\), our log discrepancy is
\[
 a(E;X,B)=1+\ord_E\bigl(K_W-p^*(K_X+B)\bigr).
\]
Log canonicity means nonnegative log discrepancies, and klt means
strictly positive log discrepancies. We use the usual dlt and slc
conventions. A \(\Q\)-Cartier divisor \(L\) is nef if \(L\cdot C\ge0\)
for every integral curve \(C\). Its numerical dimension, when nef, is
\[
 \nu(L)=\max\{j: L^jH^{\dim X-j}>0\},
\]
with \(H\) ample. For non-nef pseudo-effective divisors, statements
about numerical dimension use Nakayama's \(\kappa_\sigma\), as
specified in the relevant reduction. These notions are not
interchanged without a hypothesis that justifies doing so.

The Iitaka dimension is the maximum dimension of the images of the
complete systems of integral multiples, and is \(-\infty\) when all
these systems are empty. Numerical and rational-linear equivalence
are denoted by \(\equiv\) and \(\sim_{\Q}\), respectively. Unless a
different base field is specified, the constructions are over \(\C\).
A very general point avoids a countable union of proper closed
subvarieties. The transfer to arbitrary algebraically closed
characteristic-zero fields is made only after the complex argument
has been completed.
